% !TeX root = ../../thesis.tex
\chapter{Object Completion}
\label{ch:7_Object_Completion}

% ── New glossary entries ─────────────────────────────────────────────────────
\newglossaryentry{DiT}{name={DiT}, description={Diffusion Transformer}}
\newglossaryentry{VAE}{name={VAE}, description={Variational Auto-Encoder}}
\newglossaryentry{VQ-VAE}{name={VQ-VAE}, description={Vector-Quantised Variational Auto-Encoder}}
\newglossaryentry{COACD}{name={COACD}, description={Collision-Aware Approximate Convex Decomposition}}
\newglossaryentry{ABF}{name={ABF}, description={Angle-Based Flattening}}
\newglossaryentry{FoV}{name={FoV}, description={Field of View}}

% ─────────────────────────────────────────────────────────────────────────────
\section{Introduction}
\label{sec:7_intro}

% Paragraph purpose: situate the chapter in the pipeline and state the problem.
This chapter addresses the object completion stage of the pipeline. Each isolated foreground object delivered by Chapter~\ref{ch:5_Scene_Segmentation} must be turned into an independent, complete and textured mesh. Two outputs are needed simultaneously: a complete geometry that covers the surface the scanner could not observe, and a complete texture that covers the corresponding appearance. Both are unknown beyond the visible portion of the object, and unlike the planar walls and floors of Chapter~\ref{ch:6_Scene_Completion}, indoor furniture and equipment have arbitrary topology, free-form curvature and category-specific detail, so the per-plane reconstruction strategy used for the environment does not apply. A learned model is therefore needed to infer the missing parts.

% Paragraph purpose: introduce the two methods as a response to this problem.
Two complementary methodologies are proposed in response, summarised at a glance in Table~\ref{tab:7_method_comparison} and outlined in Figure~\ref{fig:7_intro}. The first relies on a voxelised distance field of the partial object together with the per-object occlusion grid of Chapter~\ref{ch:5_Scene_Segmentation}, and uses AutoSDF~\cite{mittal_autosdf_2022} to complete the missing voxels (Section~\ref{sec:7_voxel}). The second uses a cropped image of the object together with its partial point cloud, and uses Hunyuan3D-Omni~\cite{hunyuan3d_hunyuan3d-omni_2025} to generate the complete shape from the image (Section~\ref{sec:7_image}). The two methods sit on opposite ends of the input spectrum: one consumes a geometry-only voxel representation, the other an image-led representation with geometric guidance. Together they cover the range of failure modes that arise when one or the other input source is sparse, noisy or out of distribution.

\begin{figure}[hbht]
 \centering
 \includegraphics[width=1\textwidth]{Object_completion_intro}
 \caption[Object completion introduction]{Each isolated partial object (left) is completed into a fully textured mesh (right) using either a voxel occlusion grid (middle top) or an image view (middle bottom).}
 \label{fig:7_intro}
\end{figure}

\begin{table}[h]
    \centering
    \caption[Voxel-based vs.\ image-based object completion]{Comparison of the two object completion methods presented in this chapter along the axes that matter for the pipeline of this thesis.}
    \label{tab:7_method_comparison}
    \begin{tabularx}{\textwidth}{l >{\raggedright\arraybackslash}X >{\raggedright\arraybackslash}X}
        \toprule
        \textbf{Aspect} & \textbf{Voxel-based (AutoSDF)} & \textbf{Image-based (Hunyuan3D-Omni)} \\
        \midrule
        Main input               & Partial textured mesh + occlusion grid       & Image of the object (any type) \\
        Auxiliary input          & Refined oriented bounding box                & Partial point cloud as geometric guide \\
        Resolution               & $64^3$ voxel grid (fixed)                    & Octree up to $256^3$ at decoding \\
        Alignment sensitivity    & High; relies on Chapter~\ref{ch:5_Scene_Segmentation} OBB refinement & Lower; image provides global cues \\
        Texture detail           & Patch-quilted from observed regions          & Generated from image with delighting \\
        Preserves observed parts & Partial (\gls{VQ-VAE} alters input)          & Yes (point cloud acts as direct control) \\
        \bottomrule
    \end{tabularx}
\end{table}


% ─────────────────────────────────────────────────────────────────────────────
\section{Related Work}
\label{sec:7_related}

This section reviews the three areas most relevant to the contributions of this chapter: 3D object completion networks, texture completion for completed objects and Material detection.


% ─────────────────────────────────────────────────────────────────────────────
\subsection{3D Object Completion Networks}
\label{sec:7_rw_completion}

% Paragraph purpose: retrieval-based first (oldest paradigm, weakest assumption)
Early approaches to object completion do not predict the missing geometry, but \textit{replace} the partial scan with the closest match from a CAD database. RetrievalFuse~\cite{siddiqui_retrievalfuse_2021} learns embeddings of both scene chunks and database items so that the most similar prior can be assembled around a partial scan. Earlier hybrid pipelines~\cite{li_hybrid_2024} segment the room point cloud into structural and object components, reconstruct the structural elements from detected planes, and replace the detected objects with fitted CAD models. Database-driven completion produces clean, watertight meshes for free, but its assumption is also its weakness: any object without a close neighbour in the database is replaced by a wrong category, and the heterogeneous furniture and equipment encountered in real indoor scans are exactly where such mismatches occur. For this reason, retrieval-based methods are not considered further in the rest of this chapter.

% Paragraph purpose: point-based methods
Point-based methods operate directly on a fixed-size point cloud. Point-Voxel Diffusion~\cite{zhou_3d_2021} predicts the complete shape through a 3D diffusion process, and IF-Net~\cite{chibane_implicit_2020} uses implicit features extracted from the partial point cloud to predict occupancy at arbitrary query positions. These models recover the broad shape reliably but lose fine detail in dense regions and struggle with large missing parts, because the point representation does not encode a continuous surface and the partial sampling drops fine features along with the missing geometry.

% Paragraph purpose: voxelised distance-field methods + AutoSDF figure
Voxelised distance-field methods convert the partial geometry to an SDF or, for non-watertight scans, a UDF, both defined in Chapter~\ref{ch:background}, and complete the field on a regular voxel grid. AutoSDF~\cite{mittal_autosdf_2022} learns a discrete autoregressive prior over voxelised SDF grids using a \gls{VQ-VAE} and predicts the missing voxels conditioned on the observed ones (Figure~\ref{fig:7_autosdf}). SDFusion~\cite{cheng_sdfusion_2023} extends this with multimodal conditioning such as text or images, and DiffComplete~\cite{chu_diffcomplete_2023} reformulates the problem as a diffusion process. PatchComplete~\cite{rao_patchcomplete_2022} and ShapeFormer~\cite{yan_shapeformer_2023} use multi-resolution patch priors and a sparse vector-quantised implicit representation respectively to improve generalisation to unseen categories. XCube~\cite{ren_xcube_2023} scales the same idea to large scenes through a hierarchical sparse voxel latent diffusion model. The Interactive Voxel Editor of~\cite{wegen_interactive_2022} adds a complementary contribution: it shows that the voxel-level mask of the to-be-completed region can be edited by hand in a 3D interface, which is the workflow followed in~\cite{vermandere_guided_2025}.

\begin{figure}[hbht]
 \centering
 \includegraphics[width=1\textwidth]{Auto_SDF}
 \caption[AutoSDF method overview]{AutoSDF~\cite{mittal_autosdf_2022}: the partial SDF is encoded into discrete latent tokens by a \gls{VQ-VAE}; a transformer then autoregressively predicts the latent tokens for the voxels marked as to-be-completed, conditioned on the observed ones, and the result is decoded back to a complete SDF.}
 \label{fig:7_autosdf}
\end{figure}

% Paragraph purpose: image-based methods (no Hunyuan2.0 mention here) + figure
Image-based generative methods bypass the voxelisation step and predict 3D geometry directly from one or more images. SAM3D~\cite{team_sam_2025} combines segmentation masks with a large-scale latent flow transformer to estimate coarse geometry, pose and texture, decoded into meshes and Gaussian splats. Trellis~\cite{xiang_structured_2025} introduces a structured latent representation that can be decoded into multiple 3D formats through rectified flow transformers. Hunyuan3D-Omni~\cite{hunyuan3d_hunyuan3d-omni_2025} proposes a unified cross-modal framework that accepts an image together with optional 3D priors such as point clouds, bounding boxes or skeletons, enabling controllable generation under partial input (Figure~\ref{fig:7_hunyuan3d}).

\begin{figure}[hbht]
 \centering
 \includegraphics[width=1\textwidth]{Hunyuan_3D}
 \caption[Hunyuan3D-Omni method overview]{Hunyuan3D-Omni~\cite{hunyuan3d_hunyuan3d-omni_2025}: a DINO-based image encoder and a point encoder produce condition features that are concatenated and fed to a Diffusion Transformer; the decoded latent is converted to an SDF and meshed.}
 \label{fig:7_hunyuan3d}
\end{figure}

% Paragraph purpose: chosen methods in their own closing paragraph
Two of the above are selected for this chapter. AutoSDF is selected for the voxel-based pipeline because it explicitly supports voxel-level guidance over the to-be-completed region, which is exactly the signal that the per-object occlusion grid of Chapter~\ref{ch:5_Scene_Segmentation} provides automatically. Hunyuan3D-Omni is selected for the image-based pipeline because it accepts the partial point cloud as an additional geometric control input, which keeps the generated shape anchored to the observed surface rather than letting it drift to whatever the image alone suggests.


% ─────────────────────────────────────────────────────────────────────────────
\subsection{Texture Completion}
\label{sec:7_rw_texture}

% Paragraph purpose: pure 2D inpainting methods that need a UV map
Once the geometry is complete, the missing texture can be recovered by treating the object's texture map as an image and inpainting the unobserved regions. Classical 2D approaches model textures as Gaussian processes and perform conditional sampling for inpainting~\cite{galerne_texture_2017}, and patch-based image quilting~\cite{efros_image_2001} synthesises new texture by stitching together blocks from a reference region. Learning-based methods such as MAT~\cite{li_mat_2022} and CM-GAN~\cite{zheng_cm-gan_2022} employ transformer architectures and cascaded modulation to handle large missing regions, while Inpaint Anything~\cite{yu_inpaint_2023} combines segmentation models with diffusion-based inpainters. LaMa~\cite{suvorov_resolution-robust_2022}, used for scene completion in Chapter~\ref{ch:6_Scene_Completion}, is particularly effective at completing wide masks on repetitive structures but is less suited to the detailed and category-specific textures of foreground objects.

% Paragraph purpose: 2D inpainting needs a good UV map
A common assumption of all of these methods is that adjacent regions in the texture image correspond to adjacent regions on the surface, since the inpainting decisions are driven entirely by the local pixel neighbourhood. This assumption breaks for scanned objects, whose UV maps are typically optimised to maximise the use of texture-image space and therefore scatter neighbouring 3D faces across disconnected 2D patches~\cite{maggiordomo_texture_2021}. Several remedies have been proposed. Semantic UV mapping~\cite{vermandere_semantic_2024} re-unwraps the object in semantically meaningful chunks so that connected 3D regions remain connected in 2D. Part-aware approaches go further by first decomposing the mesh into geometrically simpler parts, for which standard unwrapping algorithms such as \gls{ABF}~\cite{maggiordomo_texture_2021} produce low-distortion charts. \gls{COACD}~\cite{wei_approximate_2022} performs this decomposition through collision-aware approximate convex decomposition and a tree search, producing nearly convex parts that each unwrap cleanly while preserving spatial proximity.

% Paragraph purpose: 3D vertex-color methods (IF-Net Texture) - state limitation
A different family avoids the UV problem entirely by predicting colour directly in 3D. IF-Net Texture~\cite{chibane_implicit_2021} extends IF-Net with per-vertex colour prediction so that the geometry and the appearance are produced together. Because the colour is stored per vertex, the achievable resolution is bounded by the vertex density of the mesh, which on scanned objects is often coarser than the texture-image resolution of UV-based methods, so fine surface detail such as small patterns or sharp colour boundaries is smoothed out (Figure~\ref{fig:7_ifnet_texture}).

\begin{figure}[hbht]
 \centering
 \includegraphics[width=1\textwidth]{IF-Net}
 \caption[IF-Net Texture method overview]{IF-Net Texture~\cite{chibane_implicit_2021}: per-vertex colour is predicted alongside the implicit surface by extending IF-Net with a colour decoder head. The geometry and appearance are produced together, but the resolution of the colour is bounded by the vertex density of the mesh. From \cite{chibane_implicit_2021}}
 \label{fig:7_ifnet_texture}
\end{figure}

% Paragraph purpose: TEXTure - also needs UV but smaller patches
TEXTure~\cite{richardson_texture_2023} sits between the two families: it generates 2D renders of the completed object through a diffusion model and projects them back onto the surface. The projection step still writes into a UV map, but because each render covers only the geometry visible from one view, the texture is composed of many small UV-aligned chunks rather than a single object-wide chart, which lifts the alignment requirement on the unwrap.

% Paragraph purpose: Hunyuan2.0 paint
Image-conditioned generative texturing such as the Hunyuan3D~2.0 paint pipeline~\cite{zhao_hunyuan3d_2025} builds on this multi-view idea and adds three ingredients: a reference-net that preserves appearance details across views, geometry conditioning that ensures the synthesised textures adhere to the surface, and multi-view attention that keeps the result consistent between viewpoints. The output is a baked texture map produced directly from a single input image.


% ─────────────────────────────────────────────────────────────────────────────
\subsection{Material Detection}
\label{sec:7_rw_material}

Most of the texture methods above treat the object as a single texture surface, but real furniture combines several distinct materials. Detecting these materials and reasoning per material rather than per object improves both segmentation boundaries and inpainting quality, because each material provides its own coherent reference region.

In 2D, Segment Anything (SAM)~\cite{kirillov_segment_2023} detects similar visual regions in an image and produces a large set of disjoint masks. Materialistic~\cite{sharma_materialistic_2023} specialises in detecting similar materials in a single image, although with limited granularity, and view-based approaches such as Material-Based Segmentation of Objects~\cite{stets_material-based_2019} segment the object from rendered views. Reflective materials remain difficult because they reflect the environment rather than themselves; Multimodal Material Segmentation~\cite{liang_multimodal_nodate} addresses this by combining RGB, near-infrared and polarised images.

In 3D, TextureNet~\cite{huang_texturenet_2018} encodes the colour of feature points into more distinctive descriptors and uses these for segmentation, which extends naturally to per-material labels within a single object. The voxel-based texture pipeline used later in this chapter follows this 3D direction: it combines an SAM-based 2D segmentation with a 3D classifier that labels every face of the completed mesh with a material index.


% ─────────────────────────────────────────────────────────────────────────────
\section{Voxel-Based Object Completion}
\label{sec:7_voxel}
 
The voxel-based methodology turns the partial textured mesh and the per-object occlusion grid of Chapter~\ref{ch:5_Scene_Segmentation} into a fully textured completed mesh through three stages (Figure~\ref{fig:7_voxel_method}): a voxelised distance field of the partial object is completed by AutoSDF, the missing material indices on the resulting mesh are predicted by IF-Net, and the per-material textures are inpainted on a re-unwrapped UV map.
 
\begin{figure}[hbht]
 \centering
 \includegraphics[width=1\textwidth]{Object_completion_voxel_method}
 \caption[Voxel-based object completion pipeline]{Overview of the voxel-based object completion pipeline. The partial point cloud (top left) is converted to a UDF inside the refined bounding box, completed by AutoSDF using the occluded voxels as the to-be-completed selection, then textured through material segmentation and per-material inpainting, resulting in a completed mesh (bottom left).}
 \label{fig:7_voxel_method}
\end{figure}
 
% Note on input: where the partial textured mesh comes from, since the scanner produces points and not meshes
The partial textured mesh expected by this pipeline is not produced directly by the scanner, since terrestrial laser scanners and mobile mapping systems output unstructured coloured point clouds rather than meshes with a UV map. Two routes are used in practice to obtain it. When multiple aligned images of the object are available within the scene, a photogrammetric reconstruction pipeline yields a textured surface directly. When only the panoramic image is available, the partial point cloud is first meshed through Poisson surface reconstruction~\cite{kazhdan_screened_2013} and the panorama is then projected onto the resulting surface to produce a texture, following the projection scheme already used in Chapter~\ref{ch:6_Scene_Completion}. Either route gives a partial mesh that covers the observed surface and carries an existing UV layout and texture image.
 
 
% ── 7.3.1 ────────────────────────────────────────────────────────────────────
\subsection{Voxel Preparation}
\label{sec:7_voxel_prep}
 
% Motivation first
AutoSDF expects two inputs: a voxelised distance field of the partial object and a binary mask of the voxels that need to be completed. The first encodes what was observed; the second tells the network where to invent geometry. Both must be produced before any completion can happen.
 
\paragraph{Voxel grid and distance field.}
The refined oriented bounding box of each object is divided into a $64^3$ voxel grid $\mathbf{G}$, matching the input resolution of AutoSDF. Because partial scans are inherently non-watertight, the partial mesh is converted to a Unsigned Distance Field $f(\mathbf{x})_\text{udf}$ rather than a standard SDF: any signed convention would require an inside-outside decision that the partial surface cannot reliably provide, as motivated in Chapter~\ref{ch:background}. The UDF is sampled at the centre of each voxel to obtain a per-voxel scalar $g_{ijk} = f(\mathbf{x}_{ijk})_\text{udf}$.
 
\paragraph{To-be-completed voxel mask.}
% Motivation: prior workflows required manual or symmetry-based mask; the occlusion grid removes this need
In earlier voxel-based completion workflows, the to-be-completed mask had to be defined either by extending the bounding box past the observed surface~\cite{mittal_autosdf_2022} or by hand-picking voxels in a 3D editor~\cite{wegen_interactive_2022, vermandere_guided_2025}. Both options require an extra step that is either coarse or labour-intensive. The per-object occlusion grid produced by the voxelised ray marching of Chapter~\ref{ch:5_Scene_Segmentation} solves this problem directly: every voxel inside the bounding box has already been classified as observed-empty, occupied or occluded, and the occluded voxels are exactly the regions where geometry could still be hidden. A voxel is therefore added to the completion set if its label is \textit{occluded} and it is not already occupied by observed surface. If the occlusion grid was computed at a different resolution from the AutoSDF grid, it is resampled to $64^3$ through nearest-neighbour upsampling so that the two grids align voxel-for-voxel. Because the occlusion grid was computed inside the oriented bounding box rather than along the world axes, the resulting mask is automatically aligned with the object's principal axes.
 
 
% ── 7.3.2 ────────────────────────────────────────────────────────────────────
\subsection{Geometry Completion}
\label{sec:7_voxel_geom}
 
% Motivation: the completion is probabilistic; multiple candidates are generated, the best one matching the observed surface is kept
With the UDF and the mask in place, AutoSDF treats the voxel grid as a sequence of tokens, autoregressively predicting the latent codes for the masked voxels conditioned on the observed ones, and decoding the result back to a complete SDF. Because the prediction is sampled rather than deterministic, the network returns a number of candidate completions per input (Figure~\ref{fig:7_voxel_candidates}). This variability is useful for objects with ambiguous geometry (for instance a partially scanned chair whose back is hidden), but it requires a selection step to commit to a single output.


\begin{figure}[hbht]
 \centering
 \includegraphics[width=1\textwidth]{Object_completion_multi_outputs}
 \caption[Multiple candidate completions from AutoSDF]{AutoSDF returns multiple candidate completions per input (top: input UDF, bottom: four candidate completions). The candidate whose decoded surface lies closest to the input UDF zero-level set is retained.}
 \label{fig:7_voxel_candidates}
\end{figure}


The candidate with the smallest mean Euclidean distance between its decoded surface and the input UDF zero-level set is retained, since this is the candidate whose observed parts have been altered the least by the \gls{VQ-VAE} encoding.
 
% Bridge to next step: SDF -> mesh
The output of the network is a completed SDF on the same $64^3$ grid. Because the next stage of the pipeline operates on explicit triangle geometry, the SDF is converted to a mesh through Marching Cubes~\cite{lorensen_marching_1998}, which extracts the zero level set as a watertight triangulated surface. This meshed result fills the previously occluded region while remaining consistent with the observed geometry at the boundary, and is the input to the texture stage.
 
 
% ── 7.3.3 ────────────────────────────────────────────────────────────────────
\subsection{Texture Completion}
\label{sec:7_voxel_texture}
 
% Motivation: the completed mesh has no UVs and no texture. Use the partial textured mesh introduced at the start of the section as the texture reference.
The completed mesh produced in Section~\ref{sec:7_voxel_geom} has no UV map and no texture: it was extracted from a scalar field rather than from a parameterised surface, so its appearance has to be reconstructed from scratch from the partial textured mesh introduced at the start of this section. The texture stage proceeds in four steps (Figure~\ref{fig:7_voxel_texture_pipeline}): a part-aware UV layout is generated for the completed mesh, the partial texture is baked onto this new layout where the two surfaces overlap, the material indices are classified on the observed region and extended to the unobserved one, and each material is inpainted independently in 2D.

\iffalse
\begin{figure}[hbht]
 \centering
 \includegraphics[width=1\textwidth]{Object_completion_voxel_texture}
 \caption[Voxel texture pipeline]{Texture completion pipeline. The completed mesh is decomposed with \gls{COACD} and re-unwrapped per part. The partial texture is baked onto the new UV layout where the completed mesh overlaps the partial one. SAM segments the baked texture into material patches that are grouped through feature similarity. IF-Net extends the resulting material indices over the unobserved surface, and patch-based image quilting inpaints each material region independently.}
 \label{fig:7_voxel_texture_pipeline}
\end{figure}
\fi
 
\paragraph{Part-aware UV unwrapping.}
% Motivation: completed mesh has no UVs; standard unwrapping gives a poor layout; COACD into parts solves this
\begin{wrapfigure}{r}{0.33\textwidth}
 \centering
 \includegraphics[width=0.33\textwidth]{Half_chair_decomposition}
 \caption[COACD convex decomposition]{Convex decomposition of a completed chair mesh with \gls{COACD}; each part is shown in a random colour.}
 \label{fig:7_coacd}
\end{wrapfigure}
A standard global UV unwrap of the completed mesh would produce a scattered layout with many small charts and arbitrary seams, which is exactly the situation that the texture stages of state-of-the-art inpainters do not handle well. To give the downstream steps a UV layout in which adjacency is meaningful, the completed mesh is first decomposed into nearly convex parts with \gls{COACD}~\cite{wei_approximate_2022}. Each part is geometrically simple enough to be unwrapped without large distortion or arbitrary cuts, so a standard \gls{ABF} unwrap~\cite{maggiordomo_texture_2021} produces a chart in which faces that are adjacent on the 3D part remain adjacent in the UV image. Combined across all parts, the result is a semantic UV layout~\cite{vermandere_semantic_2024} that preserves the spatial proximity of the original mesh, at the cost of a slightly larger texture footprint than a fully optimised packing.
 
\paragraph{Partial texture baking.}
% Motivation: the new mesh has good UVs but no texture; the partial mesh has the texture but a different (often bad) UV layout. Bake from the partial onto the new using 3D correspondence.
With a UV layout in place, the partial texture is transferred from the partial mesh onto the completed mesh through a 3D-to-3D correspondence. For each point on the completed mesh, the closest point on the partial mesh is found through a k-d tree query: when this closest point lies within a distance threshold (one voxel of $\mathbf{G}$), the colour of the partial texture at that point is written into the corresponding UV pixel of the completed mesh; otherwise the pixel is marked as unknown. Because the two meshes share their observed surface up to small local deviations, this baking populates the UV map of the completed mesh wherever the original scan covered the object, and leaves the regions reconstructed by AutoSDF empty (Figure~\ref{fig:7_voxel_texture_pipeline}). The result is a partial texture on a clean, semantically organised UV layout, which is exactly what the next steps need.
 
\paragraph{Material segmentation and 3D extension.}
% Motivation: SAM doesn't know about materials, only visual patches; we need to group patches into materials
The baked partial texture on the completed mesh is segmented into patches using SAM~\cite{kirillov_segment_2023}, which produces a large set of disjoint regions that share visual similarity. Patches that belong to the same material are then grouped through the cosine similarity of EfficientNet~\cite{tan_efficientnet_2019} features:
 
\begin{equation}
    \label{eq:7_patch_grouping}
    \begin{gathered}
    \boldsymbol{S} = \bigcup_{i} \left\{ S_i = \{P_i, P_j\} \,\Big|\, \forall P_i, P_j \in \boldsymbol{P} : \frac{f_\text{feat}(P_i) \cdot f_\text{feat}(P_j)}{\|f_\text{feat}(P_i)\| \, \|f_\text{feat}(P_j)\|} \geq t_c \right\}.
    \end{gathered}
\end{equation}
 
The similarity threshold $t_c$ is tuned so that no more than eight groups remain, since the next step encodes the material index in three binary channels and can therefore carry $2^3 = 8$ labels. The boundaries of each group are baked into the mesh as new edges so that every face carries a single material index throughout the rest of the pipeline.

\begin{figure}[hbht]
 \centering
 \includegraphics[width=1\textwidth]{Material_segmentation}
 \caption[Material segmentation pipeline]{Material segmentation: UV unwrapping of the completed mesh (left), full SAM patch segmentation on the baked partial texture (centre), and the final material grouping after EfficientNet feature clustering (right).}
 \label{fig:7_material_segmentation}
\end{figure}

% Motivation: now extend the partial labelling to the newly created faces
The faces produced by AutoSDF still need a material label, since the baking step left them empty. IF-Net~\cite{chibane_implicit_2021} is used for this prediction. Material indices are coarser than the eventual texture: they only need to identify which patch a face belongs to, not what its colour is. A per-vertex prediction is therefore sufficient, and the resolution limitation that made IF-Net unsuitable for full texture recovery does not affect material indices. The completed mesh is sampled as a coloured point cloud where each existing point carries its material index colour, voxelised at the AutoSDF resolution, and passed through the IF-Net encoder. The decoder predicts the material index at every query point, extending the observed labelling onto the previously unseen regions.

\begin{figure}[hbht]
 \centering
 \includegraphics[width=1\textwidth]{Material_prediction}
 \caption[IF-Net material index prediction]{Material index prediction: partial material indices projected onto the completed mesh (left) and the full material labelling after IF-Net extension to the unobserved geometry (right).}
 \label{fig:7_material_prediction}
\end{figure}

\paragraph{Per-material 2D inpainting.}
% Motivation: now that we have material masks, we can inpaint each material independently from a clean reference
The per-vertex material indices are projected back to the UV plane, producing a 2D material mask aligned with the baked partial texture. For each material set $S_i$, the observed patches $P \in S_i$ act as the reference texture and all other regions are masked out, so the inpainter receives a clean single-material reference instead of a mixed one. The missing texture is then synthesised through patch-based image quilting~\cite{efros_image_2001}: square blocks $\{B_1, B_2, \ldots, B_n\}$ from $S_i$ are stitched together with overlap, and the cut between adjacent blocks is chosen to minimise the cumulative neighbour-contrast error
\begin{equation}
    \label{eq:7_inpainting}
    \begin{gathered}
        E_{ij} = e_{ij} + \min(E_{i-1,j-1},\, E_{i-1,j},\, E_{i-1,j+1}),
    \end{gathered}
\end{equation}
with $e_{ij} = f(B_i, B_j)$ the boundary cost between blocks $B_i$ and $B_j$. Each material is inpainted independently, which keeps the boundaries sharp and prevents textures from one material bleeding into another. The result is a fully textured completed mesh.

\begin{figure}[hbht]
 \centering
 \includegraphics[width=1\textwidth]{Object_inpainting}
 \caption[Per-material texture inpainting]{Per-material texture inpainting: the completed mesh with the partial texture baked onto the observed geometry and material indices on the new geometry (left); the re-unwrapped UV map before inpainting (top centre) and after patch-based quilting per material (bottom centre); the final fully textured completed mesh (right).}
 \label{fig:7_object_inpainting}
\end{figure}
 

% ── 7.3.4 ────────────────────────────────────────────────────────────────────
\subsection{Experiments}
\label{sec:7_voxel_exp}

% Motivation for the evaluation choice: AutoSDF was trained on ShapeNet; standard furniture is in-distribution. Real shortcomings only become visible on categories outside that prior, so an industrial dataset is used in addition to Matterport.
The voxel-based method is evaluated on two datasets: a quantitative evaluation on a dedicated industrial dataset to expose the behaviour outside the AutoSDF training prior, and a qualitative evaluation on Matterport3D~\cite{chang_matterport3d_2017} to illustrate the typical performance on in-distribution scanned furniture.

\paragraph{Industrial dataset.}
% Introduce the dataset first, with its motivation
The industrial dataset was collected by~\cite{decrocq_voxel_2025} specifically to push AutoSDF outside its ShapeNet prior. It contains three new object classes, all from the chemical industry: storage barrels, rectangular pipe segments and round pipe segments. The real part of the dataset comprises 22 pipe segments and 9 storage barrels captured at a brewery on the M-block of the KU Leuven Rabot campus with a Leica P30 and a Leica BLK as part of a Scan-to-BIM exercise. A matching synthetic counterpart of 59 objects covering the same categories was modelled in Revit, including a wide variety of bends, junctions and orientations (Figure~\ref{fig:7_industrial_dataset}). Both sets are voxelised to $64^3$ UDFs at twelve orientations to give the network a balanced view of each object. The synthetic counterpart serves the dual purpose of providing fine-tuning data and a controlled evaluation regime in which the input is free of scan artefacts.

\iffalse
\begin{figure}[hbht]
 \centering
 \includegraphics[width=1\textwidth]{Object_completion_industrial_dataset}
 \caption[Industrial dataset]{The industrial dataset~\cite{decrocq_voxel_2025}: synthetic Revit models of storage barrels and round/rectangular pipe segments (left) and the corresponding real scans captured at the brewery on the M-block (right).}
 \label{fig:7_industrial_dataset}
\end{figure}
\fi

% Show first: results table, then explain
\begin{table}[h]
    \centering
    \caption[Voxel-based completion on industrial objects]{Voxel-based completion results on synthetic industrial objects, evaluated with the AutoSDF model after fine-tuning on the synthetic counterpart of the same dataset for 399 epochs~\cite{decrocq_voxel_2025}. Each row aggregates all completions whose to-be-predicted volume falls in the given percentage interval. IoU is reported in percent, Chamfer distance scaled by $10^2$. Values are mean $\pm$ standard deviation.}
    \label{tab:7_industrial}
    \resizebox{\textwidth}{!}{%
    \begin{tabular}{l c c c c}
        \toprule
        \textbf{To predict (\%)} & \textbf{IoU (\%) $\uparrow$} & \textbf{Chamfer ($\times 10^{-2}$) $\downarrow$} & \textbf{F-score $\uparrow$} & \textbf{Coverage $\uparrow$} \\
        \midrule
        $\phantom{0}0$--$10$  & $71.5 \pm 6.6$  & $0.34 \pm 0.10$ & $0.120 \pm 0.031$ & $0.127 \pm 0.033$ \\
        $10$--$20$            & $72.2 \pm 6.4$  & $0.35 \pm 0.13$ & $0.122 \pm 0.031$ & $0.128 \pm 0.034$ \\
        $20$--$40$            & $69.9 \pm 7.6$  & $0.41 \pm 0.22$ & $0.125 \pm 0.033$ & $0.133 \pm 0.035$ \\
        $40$--$60$            & $69.9 \pm 7.2$  & $0.48 \pm 0.31$ & $0.128 \pm 0.036$ & $0.136 \pm 0.039$ \\
        $60$--$80$            & $64.8 \pm 10.2$ & $0.64 \pm 0.48$ & $0.128 \pm 0.039$ & $0.137 \pm 0.043$ \\
        $80$--$100$           & $61.8 \pm 12.1$ & $1.00 \pm 0.82$ & $0.122 \pm 0.043$ & $0.129 \pm 0.048$ \\
        \bottomrule
    \end{tabular}%
    }
\end{table}

\paragraph{Quantitative results on synthetic industrial objects.}
On the synthetic counterpart of the dataset (Table~\ref{tab:7_industrial}), the IoU sits in the $60$--$75\,\%$ range expected for SDF-based completion methods on their own training distribution~\cite{rao_patchcomplete_2022, chu_diffcomplete_2023}, and performance degrades only when more than $70\,\%$ of the volume must be predicted from less than $30\,\%$ of observed geometry. Round objects benefit most from fine-tuning because their surfaces are smoother and more predictable than the bends and fittings of rectangular pipework, while rectangular pipes drop more sharply at the highest completeness intervals because their right-angle features are easily perturbed by the latent-code substitution.

\paragraph{Quantitative results on real industrial scans.}
The same model applied to the real scans of the same nine barrels produces fundamentally different results (Figure~\ref{fig:7_voxel_industrial}). After fine-tuning, the per-barrel IoU still ranges between $27\,\%$ and $97\,\%$, but the Chamfer distance systematically increases with respect to the partial input rather than decreasing, and the F-score and coverage both drop after every completion, indicating that the network is adding more incorrect surface than correct surface. The cause is the noise floor of the scans: reflective metal surfaces and open vessel tops scatter the laser pulse, leaving the partial UDF with internal artefacts that AutoSDF was never trained to ignore. Because the network was trained exclusively on synthetic, clean inputs, it interprets the noise as geometry and reconstructs around it, often producing recognisable shapes from its ShapeNet prior such as furniture or aircraft parts in place of the missing pipe.

\iffalse
\begin{figure}[hbht]
 \centering
 \includegraphics[width=1\textwidth]{Object_completion_voxel_industrial}
 \caption[Voxel-based completion on industrial objects]{Voxel-based completion on real scanned industrial objects~\cite{decrocq_voxel_2025}. The partial UDF (left) contains internal noise from reflective surfaces. The fine-tuned AutoSDF output (right) reconstructs around the noise rather than removing it, often hallucinating shapes from its training prior.}
 \label{fig:7_voxel_industrial}
\end{figure}
\fi

\paragraph{Qualitative results on Matterport3D.}
On in-distribution furniture from Matterport3D~\cite{chang_matterport3d_2017}, the method behaves as the synthetic numbers suggest. Figure~\ref{fig:7_voxel_matterport} shows representative completions using the scene-detection and bounding-box-refinement pipeline of Chapter~\ref{ch:5_Scene_Segmentation}. For categories well represented in the AutoSDF training prior (sofas, chairs, stools) the overall shape is recovered consistently and the material-aware texturing produces clean boundaries between materials, with patch-based quilting filling repetitive surfaces faithfully~\cite{vermandere_geometry_2025, vermandere_guided_2025}. Objects with more limited visible geometry still produce plausible shapes thanks to the symmetry-aided bounding box and the occlusion-grid voxel selection, but the texture inpainting degrades when the reference patch is small and the missing region is large, leading to visible repetition or to ambiguity between materials when SAM over-segments the partial UV layout.

\iffalse
\begin{figure}[hbht]
 \centering
 \includegraphics[width=1\textwidth]{Object_completion_voxel_matterport}
 \caption[Voxel-based object completion on Matterport3D]{Voxel-based completion examples on Matterport3D~\cite{chang_matterport3d_2017}. From left to right: the partial object in context, the partial UDF, the AutoSDF-completed mesh, and the final textured mesh after material segmentation and per-material inpainting.}
 \label{fig:7_voxel_matterport}
\end{figure}
\fi

The two evaluations together highlight the structural limitations of the voxel-based method: it requires the partial geometry to be both clean and inside the training distribution. Even when the partial input falls inside the distribution, the \gls{VQ-VAE} encoding alters the originally observed parts of the object during the latent-code substitution~\cite{mittal_autosdf_2022, vermandere_geometry_2025}. This is the motivation for the image-conditioned method developed next.


% ─────────────────────────────────────────────────────────────────────────────
\section{Image-Based Object Completion}
\label{sec:7_image}

% Motivation: a fundamentally different approach. The image carries the high-level shape and category cues; the partial point cloud anchors the result to the observed surface; the voxel-based pipeline's failure modes come from relying on geometry alone, which the image side fixes.
The image-based methodology is a fundamentally different approach to the same problem. Instead of treating the partial geometry as the primary signal and inferring the rest from a learned voxel prior, it treats an image of the object as the primary signal and uses the partial point cloud only as a geometric anchor. The image carries the category, the shape and the appearance in a single dense representation, which lets the generator reason about parts that the scan could not capture at all. The pipeline (Figure~\ref{fig:7_image_method}) has three stages: a cropped pinhole image of the object is rendered from the panorama, Hunyuan3D-Omni~\cite{hunyuan3d_hunyuan3d-omni_2025} produces the completed mesh from the image and the partial point cloud, and the Hunyuan3D~2.0 paint pipeline~\cite{zhao_hunyuan3d_2025} adds the texture from the same image.

\iffalse
\begin{figure}[hbht]
 \centering
 \includegraphics[width=1\textwidth]{Object_completion_image_method}
 \caption[Image-based object completion pipeline]{Overview of the image-based object completion pipeline. A pinhole view of the object is rendered from the panorama using the bounding-box corners as field-of-view bounds, then fed together with the partial point cloud to Hunyuan3D-Omni for geometry generation and to Hunyuan3D-Paint for texturing.}
 \label{fig:7_image_method}
\end{figure}
\fi


% ── 7.4.1 ────────────────────────────────────────────────────────────────────
\subsection{Image Cropping}
\label{sec:7_image_crop}

% Motivation, more descriptive: WHY do we crop, and HOW (pinhole re-projection rather than masking)
The panoramic image of Chapter~\ref{ch:6_Scene_Completion} captures the entire surrounding scene, but the completion network needs a view in which the target object is dominant and unrelated geometry is absent. Feeding the full panorama would force the network to also reason about furniture, walls and clutter that have no bearing on the object at hand. The bounding-box corners produced in Chapter~\ref{ch:5_Scene_Segmentation} bound the volume the object occupies, so they also bound the angular extent that the object occupies in the panorama, which is exactly the information needed to compute a focused view.

% Step 1: project corners to panorama pixel coordinates
The eight corners $\mathbf{p}_i = (x_i, y_i, z_i)^\top$ of the refined oriented bounding box are first expressed in the panorama camera frame as $\mathbf{p}_i^c = \mathbf{R}(\mathbf{p}_i - \vec{\mathbf{t}})$, using the same scanner transform $\mathbf{T} = [\mathbf{R} \mid \vec{\mathbf{t}}]$ that placed the panorama in world coordinates. The camera-frame coordinates are converted to spherical coordinates
\begin{equation}
    \label{eq:7_spherical}
    \begin{gathered}
    \theta_i = \operatorname{atan2}(x_i^c,\, z_i^c), \qquad
    \phi_i = \arcsin\!\left(\frac{y_i^c}{\|\mathbf{p}_i^c\|}\right),
    \end{gathered}
\end{equation}
and finally to equirectangular pixel coordinates for an image of width $W$ and height $H$:
\begin{equation}
    \label{eq:7_equirectangular}
    \begin{gathered}
    u_i = W \left( \frac{\theta_i + \pi}{2\pi} \right), \qquad
    v_i = H \left( \frac{\pi/2 - \phi_i}{\pi} \right).
    \end{gathered}
\end{equation}

% Motivation: pinhole reprojection rather than rectangular masking
% Step 2: convert footprint to pinhole FoV and reproject
The eight projected coordinates $\mathbf{c}_i = (u_i, v_i)$ form a polygon on the panorama; a simple axis-aligned crop of this polygon would still produce an equirectangular-distorted image, which is not how natural views are taken and not what the completion network was trained on. Instead, the angular bounds of the eight corners define the horizontal and vertical \gls{FoV} of a pinhole camera placed at the scanner position and oriented towards the centre of the object. The panorama is then resampled through this pinhole model, producing an undistorted image that frames the object as if it had been photographed in isolation by a virtual camera at the scanner location (Figure~\ref{fig:7_pinhole_crop}). This view is paired with a transparency mask derived from the projected bounding-box polygon to suppress what remains of the surrounding scene at the edges of the crop.

\iffalse
\begin{figure}[hbht]
 \centering
 \includegraphics[width=1\textwidth]{Object_completion_pinhole_crop}
 \caption[Pinhole view from the panorama]{The bounding-box corners (left, projected onto the equirectangular panorama) define the field-of-view of a virtual pinhole camera, through which the panorama is resampled to produce an undistorted close-up view of the object (right) suitable for image-based completion.}
 \label{fig:7_pinhole_crop}
\end{figure}
\fi


% ── 7.4.2 ────────────────────────────────────────────────────────────────────
\subsection{Geometry Generation}
\label{sec:7_image_geom}

% Motivation: how we use the network, not what the network is
The cropped pinhole view and the partial point cloud of the object are passed to Hunyuan3D-Omni as the image condition and the point condition respectively. The point cloud is downsampled to the network's expected input size and re-oriented to the canonical frame used during training, since the network assumes a fixed up-axis. The image is fed unaltered.

% Motivation: why the point cloud is essential here, not in the thesis as a whole
The role of the partial point cloud is to anchor the generated geometry to the observed surface. Without the point condition the network would still produce a plausible object from the image alone, but the result would not be aligned with the scanned data and could drift in scale, position or proportion. With the point condition, the generated SDF is conditioned to pass through the observed points, which plays the same role that the occluded-voxel mask played for AutoSDF: it tells the network which part of the answer is already fixed.

% Output of this step
The network outputs an SDF on an adaptive octree at up to $256^3$ resolution, from which a mesh is extracted at the zero level set through Marching Cubes~\cite{lorensen_marching_1998}. Two cleanup passes are then applied: floater removal eliminates small disconnected components that the decoder sometimes leaves outside the main shape, and degenerate-face removal collapses zero-area triangles produced at the iso-surface. The final output of this step is a watertight, untextured mesh of the complete object, aligned with the partial point cloud and ready to be textured.


% ── 7.4.3 ────────────────────────────────────────────────────────────────────
\subsection{Texture Generation}
\label{sec:7_image_texture}

% Motivation: same image, used twice
The Hunyuan3D~2.0 paint pipeline~\cite{zhao_hunyuan3d_2025} takes two inputs: the completed mesh from the previous step and the same cropped pinhole image used for geometry generation. Reusing the image keeps the texture consistent with the geometry it produced, and removes the need to find a separate appearance reference.

The pipeline proceeds in three stages. The image is first passed through a delighting module that removes baked-in highlights and shadows, producing a uniformly lit reference that is independent of the original scanner lighting. A geometry-conditioned multi-view synthesis module then renders the delit image onto the completed mesh from a series of learned camera viewpoints, using a reference-net to preserve appearance details across views, geometry conditioning to ensure that the textures adhere to the surface structure, and multi-view attention to maintain consistency between viewpoints. Finally, the generated multi-view images are baked into a single texture map, with an inpainting step that fills small uncovered regions by propagating information from nearby textured patches.

% Bring the object back to its place in the scene
The completed and textured object is generated in the canonical frame used by Hunyuan3D-Omni. To place it back in the scene at its original location, the bounding-box transformation recorded during Chapter~\ref{ch:5_Scene_Segmentation} is inverted and applied to the output mesh, returning the object to the position, orientation and scale it had in the original scan.


% ── 7.4.4 ────────────────────────────────────────────────────────────────────
\subsection{Experiments}
\label{sec:7_image_exp}

\todo{Add quantitative evaluation on V-Scan. The image-based pipeline is implemented (notebook available, generates the cropped panorama view, runs Hunyuan3D-Omni with the partial point cloud as the point control, runs Hunyuan3D-Paint on the result), but the full evaluation on V-Scan is not yet complete. Planned metrics, in line with Chapter~\ref{ch:6_Scene_Completion} and Chapter~\ref{ch:background}: Chamfer distance and F-score at 5\,cm against ground-truth object meshes from V-Scan~\cite{vermandere_synthetic_2026}, volumetric IoU at $64^3$, and image-quality metrics (PSNR, SSIM, LPIPS~\cite{zhang_unreasonable_2018}) computed on renders of the textured object from the original panorama viewpoint compared to ground-truth renders.}

\todo{Add qualitative results on Matterport3D~\cite{chang_matterport3d_2017} for the same object types shown in Figure~\ref{fig:7_voxel_matterport}, so a direct side-by-side comparison with the voxel-based method becomes possible. The notebook already produces single-object completions from cropped panorama views; results for multiple representative scenes (residential furniture, larger pieces, partially occluded objects) need to be collected and inserted here.}


% ─────────────────────────────────────────────────────────────────────────────
\section{Conclusion}
\label{sec:7_conclusion}

This chapter presented two object completion methodologies that take the partial textured meshes, the per-object occlusion grids and the panoramic image of the preceding chapters as input and produce fully textured object meshes as output.

The voxel-based method uses the per-object occlusion grid of Chapter~\ref{ch:5_Scene_Segmentation} directly as the to-be-completed voxel selection for AutoSDF~\cite{mittal_autosdf_2022, vermandere_guided_2025}, removing the manual or symmetry-based mask construction required by earlier voxel-based workflows. The texture is recovered through a part-aware UV unwrapping based on \gls{COACD}~\cite{wei_approximate_2022} and semantic UV mapping~\cite{vermandere_semantic_2024}, a SAM~\cite{kirillov_segment_2023} segmentation of the partial texture into material groups, an IF-Net~\cite{chibane_implicit_2021} extension of the material indices over the new surface, and a per-material patch-based inpainting~\cite{efros_image_2001}. Quantitative experiments on the synthetic part of the industrial dataset of~\cite{decrocq_voxel_2025} show IoU values between $61.8\,\%$ and $72.2\,\%$ across all completeness intervals, in line with state-of-the-art SDF-based completion. On the real scans of the same dataset, scan noise from reflective surfaces causes the network to reconstruct around the noise and the Chamfer distance increases after completion. These failures, together with the eight-material cap and the partial alteration of observed geometry by the \gls{VQ-VAE} encoding, motivate the second method.

The image-based method shifts the primary input from geometry to imagery. A pinhole view of the object is rendered from the panorama using the bounding-box corners as field-of-view bounds, and is passed to Hunyuan3D-Omni~\cite{hunyuan3d_hunyuan3d-omni_2025} together with the partial point cloud as a geometric anchor. The Hunyuan3D~2.0 paint pipeline~\cite{zhao_hunyuan3d_2025} then textures the result from the same image. By replacing the fixed voxel prior with a large-scale image-conditioned generator, the method generalises across a broader range of objects without per-category fine-tuning and removes the eight-material cap of the voxel-based pipeline, at the cost of a less direct link to the per-object occlusion grid.

Together with the completed environment from Chapter~\ref{ch:6_Scene_Completion}, the textured objects produced here form the static reconstruction of the indoor scene. They are passed on to Chapter~\ref{ch:8_Scene_Dynamification}, where the static meshes are turned into dynamic, part-aware and rescalable elements~\cite{vermandere_adaptive_2026}, completing the transition from a partial scan to an editable 3D environment.


%%%%%%%%%%%%%%%%%%%%%%%%%%%%%%%%%%%%%%%%%%%%%%%%%%
% Keep the following \cleardoublepage at the end of this file,
% otherwise \includeonly includes empty pages.
\cleardoublepage